\documentclass[../../../uem_mat_mei_q.tex]{subfiles}

\begin{document}
    次の空欄中\textsf{ア}，\textsf{イ}，\textsf{ウ}，\textsf{ク}に当てはまるものを選択肢の中から選び%
    その記号をマークせよ．それ以外の空欄に当てはまる0から9までの数字を解答用紙の所定の欄に%
    マークせよ．ただし，\,\myboxb{カキ}\,，%
    \,\myboxb{コサ}\:は2桁の数，\myboxb{シスセ}\:は3桁の数である．

    \vspace{.5\baselineskip}
    $a$を正の定数として

    \myspaceb%
    $\displaystyle%
        f(x)=(ax+4)^2%
    $

    \myspaceb%
    $\displaystyle%
        g(x)=(a+4)(ax^2+4)%
    $

    \myspaceb%
    $\displaystyle%
        h(x)=(ax-4)^2%
    $\\
    とおくとき，以下の問いに答えよ．

    \begin{enumerate}
        \item 放物線 $y=f(x)$ の頂点の座標は%
            $\left(\:\myboxd{ア}\,,\,\;\myboxd{イ}\:\right)$である．%
            $y=f(x)$ のグラフを$x$軸方向に\:\myboxd{ウ}\:だけ平行移動すると，$y=h(x)$ の%
            グラフに重なる．
        \item 放物線 $y=g(x)$ の頂点の座標は\\
            $\left(\:\myboxb{エ}\,,\,\;\myboxb{オ}\:a+%
            \:\myboxb{カキ}\:\right)$である．
        \item $y=f(x)$と$y=g(x)$のグラフの共有点はただ1つであり，その$x$座標は%
            \:\myboxd{ク}\:である．\\
            $y=f(x)$ と $y=h(x)$ のグラフの共有点はただ1つであり，%
            その座標は%
            $\left(\:\myboxb{ケ}\,,%
            \,\;\myboxb{コサ}\:\right)$である．%

            \vspace{.7\baselineskip}
            \textsf{ア，イ，ウ，クの解答群}
            
            {\renewcommand{\arraystretch}{1.6}%
                \vspace{-.3\baselineskip}%
                \hspace{\zw}%
                \begin{tabularx}{\dimexpr\linewidth-2\zw}{@{}LLLL@{}}
                    ⓪\; $0$ & ①\; $1$ & ②\; $-1$ & ③\; $4$ \\
                    ④\; $-4$ & ⑤\; $\myfracb{2}{a}$ &%
                    ⑥\; $-\myfracb{2}{a}$ & ⑦\; $\myfracb{4}{a}$ \\
                    ⑧\; $-\myfracb{4}{a}$ & ⑨\; $\myfracb{8}{a}$
                \end{tabularx}%
            }
        \item $y=f(x)$と$y=h(x)$のグラフと$x$軸に囲まれた図形を$A$とおく．%
            $A$の面積は $\myfracc{\myboxb{シスセ}}{\myboxb{ソ}\:a}$ である．%
            $y=f(x)$と$y=g(x)$と$y=h(x)$のグラフに囲まれた図形を$B$とおく．%
            $B$の面積は $\myfracc{\myboxb{タ}\:a}{\myboxb{チ}}$ である．%
            $A$の面積と$B$の面積が等しいときの定数$a$の値は\:\myboxb{ツ}\:である．
    \end{enumerate}
\end{document}
